\chapter{Unsupervised learning}

Up to now we only discussed \emph{supervised learning} methods, for which we have data related both to $\Xvec$ and $Y\negthinspace$. Our goal was to learn a rule to predict $Y$ from $\Xvec$. In \emph{unsupervised learning} methods, instead, we only have data on $\Xvec$. Here, the objective is to find latent structures and patterns in the data through a form of exploratory data analysis. These methods are also more subjective than supervised learning methods and can in fact lead to a supervised learning procedure.

Some examples of unsupervised learning methods are:
\begin{itemize}
    \item We can search for patterns in residuals: if we see a structure, it's probable we missed some important predictor in the model.
    \item \emph{Principal Component Regression}: we can use the principal components of $\Xvec$ as predictors in a regression model.
\end{itemize}

The two main types of unsupervised learning methods are:
\begin{enumerate}
    \item Clustering.
    \item \acf{pca}.
\end{enumerate}

\section{Clustering}

The goal of clustering is to describe the data structure via a partition into subgroups of observations. In contrast to classification, where we know the groups, here we need to deduce them from the data.

Formally, we split data $D$ into $K$ clusters, $C_1, C_2, \dots, C_K$, that form a partition, \ie:
\begin{equation}
    \bigcup_{k=1}^K C_k = D, \qquad C_i \cap C_j = \emptyset, \quad i \neq j,
\end{equation}
into meaningful (homogeneous) clusters. We can characterize homogeneity in different ways, such as:
\begin{itemize}
    \item Similarity within groups,
    \item Dissimilarity between groups.
\end{itemize}
The two main ingredients for this method are:
\begin{enumerate}
    \item The choice of an \emph{optimality criterion} to decide the best partition.
    \item Devising an \emph{optimization procedure}, in which the optimality criterion is used.
\end{enumerate}

\paragraph{Optimality criterion}

It is typically defined based on a \emph{distance measure}, \ie, observations in one cluster should be closer to each other according to this measure. A popular choice is the \emph{Euclidean distance}.

In general we want:
\begin{equation}
    C_1, \ldots, C_K = \argmin_{\substack{\Gamma_1, \ldots, \Gamma_K \\ \text{partition of } D}} \sum_{i=1}^K W(\Gamma_i),
\end{equation}
with $W(\Gamma_i)$ the sum of distances of any two objects in $\Gamma_i$. In the case of the Euclidean distance, we have:
\begin{equation}\label{eq:within_cluster_variation}
    W(\Gamma_i) = \frac{1}{2\abs{\Gamma_i}} \sum_{\xvec, \yvec \in \Gamma_i} \sum_{l = 1}^p (x_l - y_l)^2\negthinspace.
\end{equation}
This is also referred to as \emph{within-cluster variation}, since the $W$ is equivalent to:
\begin{equation}
    W(\Gamma_i) = \sum_{\xvec \in \Gamma_i} \sum_{l = 1}^p (x_l - \bar{x}_\mathit{il})^2\negthinspace, \quad \text{with } \bar{x}_\mathit{il} = \frac{1}{\abs{\Gamma_i}} \sum_{\xvec \in \Gamma_i} x_l.
\end{equation}
The $\bar{\xvec}_\mathit{il}$ is the sample mean of predictor $l$ in cluster $\Gamma_i$.

\begin{proof}
    Let's consider the case $p = 1$. Using \eqref{eq:within_cluster_variation} we have:
    \begin{equation}
        \begin{aligned}
            2W(\Gamma_i) & = \frac{1}{\abs{\Gamma_i}} \sum_{x, y \in \Gamma_i} (x - y)^2                                                                                                                                                   \\
                         & = \frac{1}{\abs{\Gamma_i}} \sum_{x, y \in \Gamma_i} (x^2 + y^2 - 2xy)                                                                                                                                           \\
                         & = \frac{1}{\abs{\Gamma_i}} \left( 2\abs{\Gamma_i} \sum_{x \in \Gamma_i} x^2 - 2\sum_{x, y \in \Gamma_i} xy \right)                                                                                              \\
                         & = 2\left(\sum_{x \in \Gamma_i} x^2 - \bar{x}_i\sum_{x \in \Gamma_i} x\right)                                                                                                                                    \\
                         & = 2\left(\sum_{x \in \Gamma_i} x^2 -2 \bar{x}_i\sum_{x \in \Gamma_i} x + \bar{x}_i\sum_{x \in \Gamma_i} x \right)                                                                                               \\
                         & = 2\left(\sum_{x \in \Gamma_i} x^2 - 2\bar{x}_i\sum_{x \in \Gamma_i} x + \abs{\Gamma_i}\bar{x}_i^2\right)              \qquad \expl{[$\textstyle\bar{x}_i\sum_{x \in \Gamma_i} x = \abs{\Gamma_i}\bar{x}_i^2$]} \\
                         & = 2\left(\sum_{x \in \Gamma_i} x^2 - \sum_{x \in \Gamma_i}2\bar{x}_i x + \sum_{x \in \Gamma_i}\bar{x}_i^2\right) = 2\sum_{x \in \Gamma_i} (x - \bar{x}_i)^2\negthinspace.
        \end{aligned}
        \vspace{-1.5em}
    \end{equation}
\end{proof}

Another formulation for the optimality criterion is the so called \emph{between-cluster variation}, that comes from the idea that minimizing the within-cluster variation is equivalent to maximizing the between-cluster variation, \ie, dissimilarity between clusters. Indeed, the overall distance between any pair of observations is a constant that can be split into the sum of the within-cluster variation and the between-cluster variation. Therefore, minimizing one is equivalent to maximizing the other.

\paragraph{Optimization algorithm}



\begin{figure}[tb]
    \centering
    \forestset{
        LT/.style = {
                delay={where content={}{shape=coordinate}{}},
                where n children=0{tier=word, baseline}{},
                for tree={
                        text height = 2ex,
                        text depth  = 0.5ex,
                        inner ysep = 0pt,
                        inner xsep = 1pt,
                        forked edge,
                        s sep = 3mm,
                    }}}

    \begin{forest}  LT
        [
        [
                [
                        [
                                [\phantom{a}]
                                    [$a$]
                            ]
                            [
                                [\phantom{a}]
                            ]
                    ]
                    [
                        [$b$]
                            [
                                [\phantom{a}]
                                    [$c$]
                            ]
                    ]
            ]
            [
                [\phantom{a}]
                    [
                        [\phantom{a}]
                            [\phantom{a}]
                    ]
            ]
        ]
        \draw[latex-latex] (-0.6, -4.1) -- (-0.6, -2.5) node[pos=0.7, right] {\scriptsize $d_2$};
        \draw[latex-latex] (-1.15, -4.1) -- (-1.15, -1.4) node[pos=0.75, right] {\scriptsize $d_1$};
        \draw[dashed] (-3, -.75) -- (3, -.75) node[pos=0.85, above] {\scriptsize $K = 2$};
        \draw[dashed] (-3, -1.75) -- (3, -1.75) node[pos=0.875, above] {\scriptsize $K = 4$};
    \end{forest}
    \caption{Example of a dendogram. $a$, $b$ and $c$ are three observations.}
    \label{fig:dendogram}
\end{figure}

The optimal problem is actually NP-hard, since there are approximately $K^n$ ways to split $n$ observations into $K$ clusters. Therefore, efficient greedy approaches are needed. However, they will only lead to local optima. Popular choices are:
\begin{description}[beginpenalty=10000]
    \item[$K$\negthinspace-means] This method uses the Euclidean distance as the optimality criterion and fixes $K$. The algorithm is the following:
          \begin{enumerate}[beginpenalty=10000]
              \item \emph{Random initialization}: assign a random cluster to each observation. This can be done \emph{uniformly} (typically used) or at random.
              \item \emph{Iterate}:
                    \begin{itemize}
                        \item For each cluster $k = 1, \ldots, K$ compute the centroid $\bar{\xvec}_k$ of the cluster.
                        \item Assign each observation $\xvec_i,\,i = 1, \ldots, n$ to the cluster whose centroid is closest to it, with respect to the Euclidean distance.
                    \end{itemize}
          \end{enumerate}

          There is no guarantee of a global maximum, so a good practice is to perform the random initialization a certain number of times and keep the best result.
    \item[Hierarchical clustering] It is a family of clustering techniques with various distance measures that can be used. Differently to $K$\negthinspace-means, $K$ is not fixed a priori. The output is a hierarchy of nested clusters, with the hierarchy represented as a \emph{dendogram} (\Cref{fig:dendogram}). In the figure, similarity is pictured vertically and not horizontally, in the sense that the observations $b$ and $c$ are ``closer'' to each other than $a$ and $b$, because the split happened at a higher level of the hierarchy (\ie, $d_1 > d_2$).

          In addition to the distance between the observations, the method needs a distance between clusters. This is called \emph{linkage criterion} and a common choice for it is the \emph{average linkage}:
          \begin{equation}
              d_{\textsc{al}}(C_i, C_j) = \frac{1}{\abs{C_i}\abs{C_j}} \sum_{\substack{\xvec \in C_i\\\yvec \in C_j}} d(\xvec, \yvec),
          \end{equation}
          where $d(\xvec, \yvec)$ is the distance between two observations, such as the Euclidean distance.

          The algorithm can proceed in a bottom-up or top-down approach:
          \begin{itemize}
              \item \emph{Agglomerative}: start with each observation in one cluster and then merge the two clusters that are closest to each other; continue until everything is in one cluster.
              \item \emph{Divisive}: start with everything in one cluster and perform recursive splitting of clusters according to the optimality criterion.
          \end{itemize}
\end{description}

\paragraph{Selection of $\bm{K}$} To select $K$, we can use the \emph{elbow method}, which consists in plotting the optimality criterion as a function of $K$ and looking for an elbow in the plot (\Cref{fig:elbow_method}). The idea is that, as $K$ increases, the optimality criterion will decrease, but after a certain point the improvement will be negligible. Therefore, we look for the point where the improvement starts to be negligible, which is exactly the elbow of the plot.

\begin{figure}[tb]
    \centering
    \begin{tikzpicture}
        \begin{axis}[
                axis lines=middle,
                xlabel={$K$},
                xlabel style={right},
                ylabel={$\displaystyle \sum_{i=1}^K W(C_i)$},
                ylabel style={
                        rotate=90,
                        yshift=.5em,
                        anchor=south,
                        at={(axis description cs:0,0.5)},
                    },
                xtick={1,2,3,4,5,6},
                ytick=\empty,
                ymin=0.1, ymax=1,
                xmin=0, xmax=7,
                width=.6\textwidth,
                height=.4\textwidth,
            ]
            \addplot[domain=1:6, thick, smooth] {0.2 + 0.15/(x-0.76)};
        \end{axis}
    \end{tikzpicture}
    \caption{Example of the curve of the elbow method for selecting $K$. In this case, the elbow is at $K = 2$.}
    \label{fig:elbow_method}
\end{figure}

\section{Clustering network data}

This is a special case of clustering, where the data are represented as a graph, with nodes and edges. It is also called \emph{community detection}. The goal is, given a network made of nodes (actors) and edges (interactions), to find unknown groupings of nodes.

We have three main approaches to this problem:
\begin{itemize}
    \item Hierarchical clustering.
    \item Spectral clustering ($K$\negthinspace-means).
    \item Stochastic block models (mixture models).
\end{itemize}

\paragraph{Optimality criterion} To define an optimality criterion we need to understand what groups have in common. Let $\ell = \tuple{C_1, \dots, C_K}$ be a partition of the nodes and $f_\mathit{ij}(\ell)$ the fraction of edges in the network that connect a node in $C_i$ to a node in $C_j$ with respect to the total number of edges. In particular, $f_\mathit{ii}(\ell)$ is the fraction of edges in cluster $C_i$. We can define the following optimality criterion (\emph{modularity}) as:
\begin{equation}
    M(\ell) = \sum_{i=1}^K (f_\mathit{ii}(\ell) - f_\mathit{ii}^*)^2\negthinspace,
\end{equation}
where $f_\mathit{ii}^*$ is the expected value of $f_\mathit{ii}(\ell)$ under a random assignment of edges to nodes. Large values of modularity suggest \emph{non-trivial} (\ie, not random) group structures.

\paragraph{Hierarchical clustering} We can apply the hierarchical clustering method, so we merge (split) nodes (clusters) that maximize the modularity.